\chapter{假设检验}

\section{假设检验的基本概念}

\subsection{假设检验问题}

在许多实际问题中，我们不仅需要对未知参数作出估计，还需要对总体的某些特性作出\textbf{判断}。例如：
\begin{itemize}
    \item 新药是否比旧药更有效？
    \item 某批产品的次品率是否超过 $3\%$？
    \item 两种测量方法的精度是否相同？
\end{itemize}

这类问题可以归结为\textbf{假设检验}（hypothesis testing）：先对总体提出一个假设，然后利用样本信息来判断这个假设是否成立。

\begin{definition}[原假设与备择假设]
将需要检验的假设称为\textbf{原假设}（零假设，null hypothesis），记为 $H_0$；
与原假设对立的假设称为\textbf{备择假设}（alternative hypothesis），记为 $H_1$。
\end{definition}

例如，检验总体均值 $\mu$ 是否等于某给定值 $\mu_0$：
\begin{itemize}
    \item 双侧检验：$H_0: \mu = \mu_0$ \quad vs.\quad $H_1: \mu \neq \mu_0$
    \item 右侧检验：$H_0: \mu \le \mu_0$ \quad vs.\quad $H_1: \mu > \mu_0$
    \item 左侧检验：$H_0: \mu \ge \mu_0$ \quad vs.\quad $H_1: \mu < \mu_0$
\end{itemize}

\subsection{两类错误}

由于决策基于样本，我们可能犯两类错误：

\begin{definition}[第一类错误与第二类错误]

\begin{itemize}
    \item \textbf{第一类错误（弃真）}：$H_0$ 为真时拒绝 $H_0$，
    \begin{equation}
        \alpha = P(\text{拒绝 } H_0 \mid H_0 \text{ 为真}).
    \end{equation}
    $\alpha$ 称为\textbf{显著性水平}（significance level），通常取 $0.01, 0.05, 0.10$。

    \item \textbf{第二类错误（取伪）}：$H_0$ 为假时接受 $H_0$，
    \begin{equation}
        \beta = P(\text{接受 } H_0 \mid H_0 \text{ 为假}).
    \end{equation}
\end{itemize}
\end{definition}

\begin{center}
\begin{tabular}{c|c|c}
    \toprule
    & $H_0$ 为真 & $H_0$ 为假 \\
    \midrule
    接受 $H_0$ & 正确决策 & 第二类错误 ($\beta$) \\
    拒绝 $H_0$ & 第一类错误 ($\alpha$) & 正确决策 \\
    \bottomrule
\end{tabular}
\end{center}

\textbf{Neyman--Pearson 原则}：在控制第一类错误概率 $\alpha$ 的前提下，使第二类错误概率 $\beta$ 尽可能小。

\begin{definition}[检验的功效]
当 $H_1$ 为真时拒绝 $H_0$ 的概率称为检验的\textbf{功效}（power）：
\begin{equation}
    \text{power} = 1 - \beta = P(\text{拒绝 } H_0 \mid H_0 \text{ 为假}).
\end{equation}
\end{definition}

\subsection{拒绝域与检验统计量}

\begin{definition}[拒绝域]
当检验统计量落入某区域 $W$ 时拒绝 $H_0$，则 $W$ 称为\textbf{拒绝域}（rejection region）或临界域。
\end{definition}

\textbf{假设检验的基本步骤：}
\begin{enumerate}
    \item 提出原假设 $H_0$ 和备择假设 $H_1$；
    \item 选取适当的检验统计量，并在 $H_0$ 为真下确定其分布；
    \item 给定显著性水平 $\alpha$，确定拒绝域 $W$；
    \item 由样本计算检验统计量的观测值，若落入 $W$ 则拒绝 $H_0$，否则不拒绝 $H_0$。
\end{enumerate}

% ============================================================
\section{正态总体均值的假设检验}

\subsection{单个正态总体均值的检验}

设 $X_1, \dots, X_n \stackrel{\mathrm{i.i.d.}}{\sim} \mathcal{N}(\mu, \sigma^2)$，检验 $H_0: \mu = \mu_0$。

\medskip
\textbf{情形 1：$\sigma^2$ 已知 — $Z$ 检验}

检验统计量：
\begin{equation}
    Z = \frac{\bar{X} - \mu_0}{\sigma/\sqrt{n}} \sim \mathcal{N}(0,1) \quad (\text{在 } H_0 \text{ 下}).
\end{equation}

\begin{center}
\begin{tabular}{c|c}
    \toprule
    备择假设 & 拒绝域 \\
    \midrule
    $H_1: \mu \neq \mu_0$ & $|Z| \ge z_{\alpha/2}$ \\
    $H_1: \mu > \mu_0$    & $Z \ge z_{\alpha}$ \\
    $H_1: \mu < \mu_0$    & $Z \le -z_{\alpha}$ \\
    \bottomrule
\end{tabular}
\end{center}

\medskip
\textbf{情形 2：$\sigma^2$ 未知 — $t$ 检验}

用样本标准差 $S$ 代替 $\sigma$：
\begin{equation}
    T = \frac{\bar{X} - \mu_0}{S/\sqrt{n}} \sim t(n-1) \quad (\text{在 } H_0 \text{ 下}).
\end{equation}

\begin{center}
\begin{tabular}{c|c}
    \toprule
    备择假设 & 拒绝域 \\
    \midrule
    $H_1: \mu \neq \mu_0$ & $|T| \ge t_{\alpha/2}(n-1)$ \\
    $H_1: \mu > \mu_0$    & $T \ge t_{\alpha}(n-1)$ \\
    $H_1: \mu < \mu_0$    & $T \le -t_{\alpha}(n-1)$ \\
    \bottomrule
\end{tabular}
\end{center}

\subsection{两个正态总体均值差的检验}

设 $X_1, \dots, X_{n_1} \stackrel{\mathrm{i.i.d.}}{\sim} \mathcal{N}(\mu_1, \sigma_1^2)$，$Y_1, \dots, Y_{n_2} \stackrel{\mathrm{i.i.d.}}{\sim} \mathcal{N}(\mu_2, \sigma_2^2)$，两样本独立。检验 $H_0: \mu_1 - \mu_2 = \delta$（常取 $\delta = 0$）。

\medskip
\textbf{情形 1：$\sigma_1^2, \sigma_2^2$ 已知}

\begin{equation}
    Z = \frac{(\bar{X} - \bar{Y}) - \delta}{\sqrt{\frac{\sigma_1^2}{n_1} + \frac{\sigma_2^2}{n_2}}} \sim \mathcal{N}(0,1).
\end{equation}

\medskip
\textbf{情形 2：$\sigma_1^2 = \sigma_2^2 = \sigma^2$ 但未知}

合并方差 $S_w^2 = \dfrac{(n_1-1)S_1^2 + (n_2-1)S_2^2}{n_1 + n_2 - 2}$，检验统计量：
\begin{equation}
    T = \frac{(\bar{X} - \bar{Y}) - \delta}{S_w \sqrt{\frac{1}{n_1} + \frac{1}{n_2}}} \sim t(n_1 + n_2 - 2).
\end{equation}

\medskip
\textbf{情形 3：$\sigma_1^2 \neq \sigma_2^2$ 且未知（Welch $t$ 检验）}

\begin{equation}
    T = \frac{(\bar{X} - \bar{Y}) - \delta}{\sqrt{\frac{S_1^2}{n_1} + \frac{S_2^2}{n_2}}},
\end{equation}
其自由度由 Welch--Satterthwaite 公式近似：
\begin{equation}
    \nu = \frac{\left(\frac{S_1^2}{n_1} + \frac{S_2^2}{n_2}\right)^2}
                 {\frac{(S_1^2/n_1)^2}{n_1-1} + \frac{(S_2^2/n_2)^2}{n_2-1}}.
\end{equation}

\subsection{配对数据的 $t$ 检验}

当两个样本不独立而是\textbf{配对}出现时（如同一组受试者在处理前后的测量值），令
\begin{equation}
    D_i = X_i - Y_i,\quad i = 1,\dots,n.
\end{equation}
假设 $D_i \stackrel{\mathrm{i.i.d.}}{\sim} \mathcal{N}(\mu_D, \sigma_D^2)$，检验 $H_0: \mu_D = 0$，使用单样本 $t$ 检验：
\begin{equation}
    T = \frac{\bar{D}}{S_D / \sqrt{n}} \sim t(n-1).
\end{equation}

% ============================================================
\section{正态总体方差的假设检验}

\subsection{单个正态总体方差的检验 — $\chi^2$ 检验}

设 $X_1, \dots, X_n \stackrel{\mathrm{i.i.d.}}{\sim} \mathcal{N}(\mu, \sigma^2)$，检验 $H_0: \sigma^2 = \sigma_0^2$。

检验统计量：
\begin{equation}
    \chi^2 = \frac{(n-1)S^2}{\sigma_0^2} \sim \chi^2(n-1) \quad (\text{在 } H_0 \text{ 下}).
\end{equation}

\begin{center}
\begin{tabular}{c|c}
    \toprule
    备择假设 & 拒绝域 \\
    \midrule
    $H_1: \sigma^2 \neq \sigma_0^2$ & $\chi^2 \le \chi^2_{1-\alpha/2}(n-1)$ 或 $\chi^2 \ge \chi^2_{\alpha/2}(n-1)$ \\
    $H_1: \sigma^2 > \sigma_0^2$    & $\chi^2 \ge \chi^2_{\alpha}(n-1)$ \\
    $H_1: \sigma^2 < \sigma_0^2$    & $\chi^2 \le \chi^2_{1-\alpha}(n-1)$ \\
    \bottomrule
\end{tabular}
\end{center}

\subsection{两个正态总体方差比的检验 — $F$ 检验}

设两样本独立，检验 $H_0: \sigma_1^2 = \sigma_2^2$（即 $\sigma_1^2 / \sigma_2^2 = 1$）。

检验统计量：
\begin{equation}
    F = \frac{S_1^2}{S_2^2} \sim F(n_1-1,\, n_2-1) \quad (\text{在 } H_0 \text{ 下}).
\end{equation}

\begin{center}
\begin{tabular}{c|c}
    \toprule
    备择假设 & 拒绝域 \\
    \midrule
    $H_1: \sigma_1^2 \neq \sigma_2^2$ & $F \le F_{1-\alpha/2}(n_1-1,\,n_2-1)$ 或 $F \ge F_{\alpha/2}(n_1-1,\,n_2-1)$ \\
    $H_1: \sigma_1^2 > \sigma_2^2$    & $F \ge F_{\alpha}(n_1-1,\,n_2-1)$ \\
    $H_1: \sigma_1^2 < \sigma_2^2$    & $F \le F_{1-\alpha}(n_1-1,\,n_2-1)$ \\
    \bottomrule
\end{tabular}
\end{center}

\begin{remark}
实际计算时常用 $F = \max(S_1^2, S_2^2) / \min(S_1^2, S_2^2)$ 简化为上侧检验。
\end{remark}

% ============================================================
\section{置信区间与假设检验之间的关系}

置信区间与假设检验之间存在深刻的\textbf{对偶关系}（duality）。

\begin{theorem}[置信区间与假设检验的对偶性]
参数 $\theta$ 的 $1-\alpha$ 置信区间恰好由所有使得双侧检验 $H_0: \theta = \theta_0$ 在显著性水平 $\alpha$ 下不被拒绝的 $\theta_0$ 组成。
\end{theorem}

换句话说：
\begin{itemize}
    \item 若 $\theta_0$ 落在 $1-\alpha$ 置信区间内 $\iff$ 不拒绝 $H_0: \theta = \theta_0$（显著性水平 $\alpha$）；
    \item 若 $\theta_0$ 不在 $1-\alpha$ 置信区间内 $\iff$ 拒绝 $H_0: \theta = \theta_0$。
\end{itemize}

\begin{example}[均值检验与置信区间的对应]
已知方差时，$\mu$ 的 $95\%$ 置信区间为 $\bar{X} \pm z_{0.025} \cdot \sigma/\sqrt{n}$。
双侧 $Z$ 检验在 $|Z| \ge z_{0.025}$ 时拒绝 $H_0: \mu = \mu_0$，
这等价于 $\mu_0$ 不在上述置信区间内。
\end{example}

\textbf{实际意义：}
\begin{itemize}
    \item 置信区间提供了\textbf{更多信息}：不仅告诉我们是否拒绝 $H_0$，还给出了参数的可能取值范围；
    \item 假设检验给出明确的\textbf{二值决策}：在需要做出是否采取行动的场合更有用。
\end{itemize}

% ============================================================
\section{样本容量的选取}

在实际问题中，我们需要在抽样前确定所需的样本容量 $n$，以保证检验达到指定的功效。

\subsection{$Z$ 检验的样本容量}

考虑 $H_0: \mu = \mu_0$ vs.\ $H_1: \mu = \mu_1$（$\sigma^2$ 已知）。

对于显著性水平 $\alpha$ 和功效 $1-\beta$，所需样本容量为：

\textbf{双侧检验：}
\begin{equation}
    \boxed{n \approx \frac{(z_{\alpha/2} + z_{\beta})^2 \sigma^2}{(\mu_1 - \mu_0)^2}}.
\end{equation}

\textbf{单侧检验：}
\begin{equation}
    \boxed{n \approx \frac{(z_{\alpha} + z_{\beta})^2 \sigma^2}{(\mu_1 - \mu_0)^2}}.
\end{equation}

其中 $\Delta = |\mu_1 - \mu_0|$ 是我们希望检测出的\textbf{最小差异}（effect size）。

\begin{example}
某工厂想检验产品重量的均值是否为 500 g。已知 $\sigma = 5$ g，
若真实的均值为 502 g，希望在 $\alpha = 0.05$，power $= 0.90$ 下能检测出来，
求所需样本容量。

由 $z_{0.025} = 1.96$，$z_{0.10} = 1.28$：
\begin{equation}
    n \approx \frac{(1.96 + 1.28)^2 \times 5^2}{(502-500)^2}
       = \frac{(3.24)^2 \times 25}{4} \approx 65.6,
\end{equation}
故取 $n = 66$。
\end{example}

\subsection{$t$ 检验的样本容量}

当 $\sigma^2$ 未知时，使用迭代方法或查 OC 曲线（operating characteristic curve）来确定样本容量。OC 曲线给出了不同样本量下第二类错误概率 $\beta$ 与参数真值的关系。

\subsection{影响样本容量的因素}

\begin{itemize}
    \item \textbf{效应量} $\Delta$：希望检测的差异越小，需要的 $n$ 越大；
    \item \textbf{显著性水平} $\alpha$：$\alpha$ 越小，需要的 $n$ 越大；
    \item \textbf{功效} $1-\beta$：功效要求越高，需要的 $n$ 越大；
    \item \textbf{总体方差} $\sigma^2$：方差越大，需要的 $n$ 越大。
\end{itemize}

% ============================================================
\section{分布拟合检验}

在许多实际问题中，总体的分布类型是未知的。我们需要根据样本检验总体是否服从某个指定的分布。

\subsection{$\chi^2$ 拟合优度检验}

\textbf{思想：} 将样本分组，比较各组的\textbf{观测频数}与在 $H_0$ 下的\textbf{期望频数}。

\begin{definition}[Pearson $\chi^2$ 统计量]
设将总体取值范围分为 $k$ 个互不相交的区间，第 $i$ 个区间的观测频数为 $f_i$，
在 $H_0$（总体服从某特定分布）下的期望频数为 $\hat{e}_i = n \cdot \hat{p}_i$，其中 $\hat{p}_i$ 为 $H_0$ 下落入第 $i$ 区间的概率。

Pearson $\chi^2$ 统计量：
\begin{equation}
    \chi^2 = \sum_{i=1}^{k} \frac{(f_i - \hat{e}_i)^2}{\hat{e}_i}.
\end{equation}
\end{definition}

当 $n \to \infty$ 时，若 $H_0$ 为真，
\begin{equation}
    \chi^2 \xrightarrow{d} \chi^2(k - r - 1),
\end{equation}
其中 $r$ 为 $H_0$ 分布中待估参数的个数。

拒绝域：$\chi^2 \ge \chi^2_{\alpha}(k - r - 1)$。

\textbf{注意事项：}
\begin{itemize}
    \item 每个区间的期望频数应 $\ge 5$，否则需合并区间；
    \item 适用于大样本（$n \ge 50$）且既可用于离散分布也可用于连续分布。
\end{itemize}

\begin{example}[检验骰子是否均匀]
掷一枚骰子 120 次，记录各面出现次数。检验 $H_0$：骰子是均匀的（$p_i = 1/6$）。
期望频数均为 $120/6 = 20$。

若观测频数为 $f = (18, 22, 17, 21, 19, 23)$：
\begin{equation}
    \chi^2 = \frac{(18-20)^2}{20} + \frac{(22-20)^2}{20} + \cdots + \frac{(23-20)^2}{20}
           = 1.7 < \chi^2_{0.05}(5) = 11.07,
\end{equation}
故不拒绝 $H_0$，认为骰子是均匀的。
\end{example}

\subsection{正态性检验概述}

判断数据是否来自正态总体是许多统计方法的前提。

\textbf{常用方法：}
\begin{itemize}
    \item \textbf{Q-Q 图（Quantile-Quantile Plot）}：将样本分位数与正态理论分位数作图，若近似为直线则支持正态性假设；
    \item \textbf{Shapiro--Wilk 检验}：适用于小样本（$n \le 50$），功效较高；
    \item \textbf{Kolmogorov--Smirnov 检验}：基于经验分布函数与理论分布函数的最大差异；
    \item \textbf{Jarque--Bera 检验}：基于样本偏度和峰度，适用于大样本。
\end{itemize}

% ============================================================
\section{秩和检验}

前面讨论的 $t$ 检验等均假设总体服从正态分布。当正态性假设不成立或数据为\textbf{定序数据}时，需使用\textbf{非参数方法}。秩和检验是最常用的非参数检验之一。

\subsection{Wilcoxon 秩和检验（两独立样本）}

\textbf{目的：} 比较两个独立总体的\textbf{位置参数}是否相同，是两样本 $t$ 检验的非参数替代。

\begin{remark}[与 $t$ 检验的对应关系]
注意区分两种不同的一一对应：
\begin{center}
\begin{tabular}{c|c}
    \toprule
    参数方法 & 非参数替代 \\
    \midrule
    单样本 $t$ 检验（$H_0: \mu = \mu_0$） & Wilcoxon \textbf{符号秩}检验 \\
    两样本 $t$ 检验（$H_0: \mu_1 = \mu_2$） & Wilcoxon \textbf{秩和}检验 \\
    \bottomrule
\end{tabular}
\end{center}
两样本 $t$ 检验是一次性直接比较两个组的均值差，用 $T = \frac{\bar{X}_1 - \bar{X}_2}{S_w\sqrt{1/n_1 + 1/n_2}}$ 作为统计量，
而不是分别对每组各做一次单样本 $t$ 检验。秩和检验也是如此——它将两组数据\textbf{合并排序}后统一比较。
\end{remark}

\medskip
\textbf{步骤：}
\begin{enumerate}
    \item 将两组数据（样本量 $n_1, n_2$）合并，从小到大排序并赋予\textbf{秩}（rank），相同值赋予平均秩；
    \item 计算第一组样本的秩和 $R_1$；
    \item 检验统计量：
    \begin{equation}
        U = n_1 n_2 + \frac{n_1(n_1+1)}{2} - R_1.
    \end{equation}
    或等价地使用 $R_1$ 本身。
\end{enumerate}

当 $n_1, n_2$ 较大时，$R_1$ 近似服从正态分布：
\begin{equation}
    \E[R_1] = \frac{n_1(n_1+n_2+1)}{2},\quad
    \Var(R_1) = \frac{n_1 n_2 (n_1+n_2+1)}{12}.
\end{equation}

\begin{example}[Wilcoxon 秩和检验完整演示]
某研究比较两种教学法对学生成绩的影响。随机抽取 A 班 6 名学生、B 班 7 名学生，
期末考试成绩如下（满分 100）：\par
\medskip
\begin{center}
\begin{tabular}{c|ccccccc}
    \toprule
    A 班 & 72 & 68 & 85 & 91 & 55 & 77 & \\
    \midrule
    B 班 & 60 & 75 & 82 & 63 & 70 & 58 & 79 \\
    \bottomrule
\end{tabular}
\end{center}

检验 $H_0$：两种教学法的成绩分布相同 vs.\ $H_1$：两种教学法成绩分布不同（$\alpha = 0.05$，双侧）。

\textbf{步骤 1：合并排序并赋予秩}

将两班共 $n_1+n_2 = 13$ 个数据从小到大排列，同分取平均秩：

\begin{center}
\begin{tabular}{c|c|c|c}
    \toprule
    原始值 & 所属班 & 排序位次 & 秩 \\
    \midrule
    55 & A & 1  & 1 \\
    58 & B & 2  & 2 \\
    60 & B & 3  & 3 \\
    63 & B & 4  & 4 \\
    68 & A & 5  & 5 \\
    70 & B & 6  & 6 \\
    72 & A & 7  & 7 \\
    75 & B & 8  & 8 \\
    77 & A & 9  & 9 \\
    79 & B & 10 & 10 \\
    82 & B & 11 & 11 \\
    85 & A & 12 & 12 \\
    91 & A & 13 & 13 \\
    \bottomrule
\end{tabular}
\end{center}

\textbf{步骤 2：计算 A 班的秩和 $R_1$}
\begin{equation}
    R_1 = 1 + 5 + 7 + 9 + 12 + 13 = 47.
\end{equation}

\textbf{步骤 3：计算检验统计量}

取 $n_1 = 6$（A 班），$n_2 = 7$（B 班）。

Mann--Whitney $U$ 统计量：
\begin{equation}
    U = n_1 n_2 + \frac{n_1(n_1+1)}{2} - R_1
      = 6 \times 7 + \frac{6 \times 7}{2} - 47
      = 42 + 21 - 47 = 16.
\end{equation}

$R_1$ 在 $H_0$ 下的期望与方差：
\begin{align}
    \E[R_1] &= \frac{n_1(n_1+n_2+1)}{2}
             = \frac{6 \times 14}{2} = 42, \\
    \Var(R_1) &= \frac{n_1 n_2 (n_1+n_2+1)}{12}
               = \frac{6 \times 7 \times 14}{12} = 49.
\end{align}

标准化（大样本近似）：
\begin{equation}
    Z = \frac{R_1 - \E[R_1]}{\sqrt{\Var(R_1)}}
      = \frac{47 - 42}{\sqrt{49}}
      = \frac{5}{7} \approx 0.714.
\end{equation}

\textbf{步骤 4：做出决策}

双侧检验，$\alpha = 0.05$，$z_{0.025} = 1.96$。

$|Z| = 0.714 < 1.96$，不拒绝 $H_0$。

也可用精确检验：查 Wilcoxon 秩和检验临界值表，对于 $n_1=6, n_2=7$（双侧 $\alpha=0.05$），
临界下界 $R_L = 30$，临界上界 $R_U = 54$。$R_1 = 47$ 落在 $[30, 54]$ 内，同样不拒绝 $H_0$。

\textbf{结论：} 在 $\alpha = 0.05$ 水平下，没有足够证据认为两种教学法的成绩分布存在显著差异。

\textbf{补充 —— $p$ 值计算：}
\begin{equation}
    p = 2 \times P(Z \ge 0.714 \mid H_0)
      = 2 \times (1 - \Phi(0.714))
      \approx 2 \times 0.2375 = 0.475.
\end{equation}
$p = 0.475 > 0.05$，与上述结论一致。
\end{example}

\textbf{特点：} 不依赖总体分布的具体形式，对异常值稳健（robust），但功效略低于 $t$ 检验（当正态假设成立时）。

\subsection{Wilcoxon 符号秩检验（配对样本）}

\textbf{目的：} 配对样本 $t$ 检验的非参数替代。与秩和检验不同，符号秩检验对应的是\textbf{配对设计}（同一受试者前后测量），
等价于对差值 $D_i$ 做单样本检验，即 $H_0$：差值的分布关于 0 对称。

\textbf{步骤：}
\begin{enumerate}
    \item 计算每对数据的差值 $D_i = X_i - Y_i$，去掉 $D_i = 0$ 的对；
    \item 对 $|D_i|$ 从小到大赋予秩；
    \item 分别计算正差值对应的秩和 $R^+$ 与负差值对应的秩和 $R^-$；
    \item 检验统计量可取 $R^+$（或 $R^-$），其分布在 $H_0$ 下可查 Wilcoxon 符号秩检验表。
\end{enumerate}

当有效对数 $n$ 较大时（通常 $n \ge 20$），$R^+$ 近似服从正态分布：
\begin{equation}
    \E[R^+] = \frac{n(n+1)}{4},\quad
    \Var(R^+) = \frac{n(n+1)(2n+1)}{24}.
\end{equation}

\begin{example}[Wilcoxon 符号秩检验完整演示]
某研究想测试服用某种补充剂是否能改善睡眠时长。招募 10 名受试者，
记录服用前与服用一个月后的睡眠时长（小时），数据如下：\par
\medskip
\begin{center}
\begin{tabular}{c|cccccccccc}
    \toprule
    受试者编号 & 1 & 2 & 3 & 4 & 5 & 6 & 7 & 8 & 9 & 10 \\
    \midrule
    服用前 ($Y_i$) & 6.2 & 7.0 & 5.8 & 6.5 & 7.5 & 5.5 & 6.0 & 7.2 & 6.8 & 5.9 \\
    服用后 ($X_i$) & 6.8 & 7.3 & 6.1 & 7.0 & 7.6 & 5.4 & 6.5 & 7.5 & 7.1 & 6.4 \\
    \bottomrule
\end{tabular}
\end{center}

检验 $H_0$：服用前后睡眠时长无差异 vs.\ $H_1$：服用后睡眠时长增加（$\alpha = 0.05$，右侧检验）。

\textbf{步骤 1：计算差值 $D_i = X_i - Y_i$ 并取绝对值}

\begin{center}
\begin{tabular}{c|c|c|c}
    \toprule
    受试者 & $X_i$（后） & $Y_i$（前） & $D_i = X_i - Y_i$ \\
    \midrule
    1  & 6.8 & 6.2 & $+0.6$ \\
    2  & 7.3 & 7.0 & $+0.3$ \\
    3  & 6.1 & 5.8 & $+0.3$ \\
    4  & 7.0 & 6.5 & $+0.5$ \\
    5  & 7.6 & 7.5 & $+0.1$ \\
    6  & 5.4 & 5.5 & $-0.1$ \\
    7  & 6.5 & 6.0 & $+0.5$ \\
    8  & 7.5 & 7.2 & $+0.3$ \\
    9  & 7.1 & 6.8 & $+0.3$ \\
    10 & 6.4 & 5.9 & $+0.5$ \\
    \bottomrule
\end{tabular}
\end{center}

没有 $D_i = 0$ 的对，有效对数 $n = 10$。

\textbf{步骤 2：对 $|D_i|$ 从小到大赋予秩}

\begin{center}
\begin{tabular}{c|c|c|c|c}
    \toprule
    受试者 & $D_i$ & $|D_i|$ & $|D_i|$ 排序 & 秩 \\
    \midrule
    5  & $+0.1$ & 0.1 & 1 & 1 \\
    6  & $-0.1$ & 0.1 & 2 & 2 \\
    2  & $+0.3$ & 0.3 & 3 & 4 \\
    3  & $+0.3$ & 0.3 & 4 & 4 \\
    8  & $+0.3$ & 0.3 & 5 & 4 \\
    9  & $+0.3$ & 0.3 & 6 & 4 \\
    4  & $+0.5$ & 0.5 & 7 & 8 \\
    7  & $+0.5$ & 0.5 & 8 & 8 \\
    10 & $+0.5$ & 0.5 & 9 & 8 \\
    1  & $+0.6$ & 0.6 & 10 & 10 \\
    \bottomrule
\end{tabular}
\end{center}

\textbf{注意结的处理：}
\begin{itemize}
    \item $|D_i| = 0.1$ 出现 2 次（位次 1, 2），秩各为 $\frac{1+2}{2} = 1.5$；
    \item $|D_i| = 0.3$ 出现 4 次（位次 3--6），秩各为 $\frac{3+4+5+6}{4} = 4.5$；
    \item $|D_i| = 0.5$ 出现 3 次（位次 7--9），秩各为 $\frac{7+8+9}{3} = 8$；
    \item $|D_i| = 0.6$ 出现 1 次，秩为 10。
\end{itemize}

带符号的秩：

\begin{center}
\begin{tabular}{c|c|c|c}
    \toprule
    受试者 & $D_i$ & 符号 & 带符号秩 \\
    \midrule
    1  & $+0.6$ & $+$ & $+10$ \\
    2  & $+0.3$ & $+$ & $+4.5$ \\
    3  & $+0.3$ & $+$ & $+4.5$ \\
    4  & $+0.5$ & $+$ & $+8$ \\
    5  & $+0.1$ & $+$ & $+1.5$ \\
    6  & $-0.1$ & $-$ & $-1.5$ \\
    7  & $+0.5$ & $+$ & $+8$ \\
    8  & $+0.3$ & $+$ & $+4.5$ \\
    9  & $+0.3$ & $+$ & $+4.5$ \\
    10 & $+0.5$ & $+$ & $+8$ \\
    \bottomrule
\end{tabular}
\end{center}

\textbf{步骤 3：计算正秩和 $R^+$ 与负秩和 $R^-$}

\begin{align}
    R^+ &= 10 + 4.5 + 4.5 + 8 + 1.5 + 8 + 4.5 + 4.5 + 8
         = 53.5, \\
    R^- &= 1.5.
\end{align}

验证：$R^+ + R^- = 53.5 + 1.5 = 55 = \frac{10 \times 11}{2} = \frac{n(n+1)}{2}$。\checkmark

\textbf{步骤 4：使用正态近似计算 $Z$ 统计量}

\begin{align}
    \E[R^+] &= \frac{n(n+1)}{4} = \frac{10 \times 11}{4} = 27.5, \\
    \Var(R^+) &= \frac{n(n+1)(2n+1)}{24}
               = \frac{10 \times 11 \times 21}{24}
               = \frac{2310}{24} = 96.25.
\end{align}

\begin{equation}
    Z = \frac{R^+ - \E[R^+]}{\sqrt{\Var(R^+)}}
      = \frac{53.5 - 27.5}{\sqrt{96.25}}
      = \frac{26}{9.811} \approx 2.650.
\end{equation}

\textbf{步骤 5：做出决策}

右侧检验，$\alpha = 0.05$，$z_{0.05} = 1.645$。

$Z = 2.650 > 1.645$，拒绝 $H_0$。

也可查 Wilcoxon 符号秩检验临界值表：$n = 10$，右侧 $\alpha = 0.05$ 时临界值为 $T_{0.05} = 44$。
$R^+ = 53.5 > 44$，同样拒绝 $H_0$。

\textbf{结论：} 在 $\alpha = 0.05$ 水平下，有显著证据表明服用该补充剂后睡眠时长有所增加。

\textbf{补充 —— $p$ 值计算（正态近似）：}
\begin{equation}
    p = P(Z \ge 2.650 \mid H_0)
      = 1 - \Phi(2.650)
      \approx 1 - 0.9960 = 0.004.
\end{equation}
$p = 0.004 < 0.01$，属于强证据拒绝 $H_0$。
\end{example}

\begin{remark}
非参数方法的优势在于\textbf{无需分布假设}，适用于小样本、有异常值或定序数据的情形。
但其代价是当分布假设实际成立时，功效略低于对应的参数方法。
\end{remark}

% ============================================================
\section{假设检验问题的 $p$ 值检验法}

\subsection{$p$ 值的定义}

传统检验方法需预先指定显著性水平 $\alpha$，然后给出"拒绝"或"不拒绝"的二值结论。但在实际应用中（尤其是在科学研究和机器学习中），更常用的是\textbf{$p$ 值}。

\begin{definition}[$p$ 值]
在原假设 $H_0$ 为真的条件下，观察到比当前样本\textbf{更极端}结果的概率。
\end{definition}

设检验统计量的观测值为 $t_{\text{obs}}$，则：
\begin{itemize}
    \item 双侧检验：$p = 2 \cdot P(T \ge |t_{\text{obs}}| \mid H_0)$；
    \item 右侧检验：$p = P(T \ge t_{\text{obs}} \mid H_0)$；
    \item 左侧检验：$p = P(T \le t_{\text{obs}} \mid H_0)$。
\end{itemize}

\subsection{$p$ 值的解释}

\begin{center}
\begin{tabular}{c|l}
    \toprule
    $p$ 值范围 & 解释 \\
    \midrule
    $p \le 0.01$ & 强证据拒绝 $H_0$ \\
    $0.01 < p \le 0.05$ & 中等证据拒绝 $H_0$ \\
    $0.05 < p \le 0.10$ & 弱证据拒绝 $H_0$ \\
    $p > 0.10$ & 不足以拒绝 $H_0$ \\
    \bottomrule
\end{tabular}
\end{center}

\textbf{$p$ 值的优势：}
\begin{itemize}
    \item 提供了比"拒绝/不拒绝"更丰富的信息——反映了\textbf{证据的强度}；
    \item 读者可根据自己的 $\alpha$ 标准自行判断；
    \item 是现代科学研究（尤其是生命科学、社会科学）中报告检验结果的标准方式。
\end{itemize}

\begin{remark}[关于 $p$ 值的常见误解]
\begin{itemize}
    \item $p$ 值不是 $H_0$ 为真的概率；
    \item $p$ 值不是效应的大小（effect size 需单独报告）；
    \item 大样本下微小的效应也可能得到很小的 $p$ 值，"显著"不等于"重要"。
\end{itemize}
\end{remark}

\subsection{$p$ 值与置信区间的关系}

$p < \alpha$ 等价于 $1-\alpha$ 置信区间不包含 $H_0$ 下的参数值。因此 $p$ 值检验法与置信区间法在数学上等价。

\begin{example}[$p$ 值检验法完整演示]
某灌装生产线标称每瓶装 500 ml。质检员随机抽取 $n = 25$ 瓶，测得
$\bar{X} = 503.2$ ml。已知灌装量服从正态分布，且长期记录表明 $\sigma = 8$ ml。
问：是否有证据表明灌装量偏离了 500 ml？（取 $\alpha = 0.05$）

\medskip
\textbf{步骤 1：建立假设}

\begin{equation}
    H_0: \mu = 500 \quad \text{vs.} \quad H_1: \mu \neq 500 \quad (\text{双侧检验}).
\end{equation}

\textbf{步骤 2：计算检验统计量}

方差已知，使用 $Z$ 检验：
\begin{equation}
    Z_{\text{obs}} = \frac{\bar{X} - \mu_0}{\sigma/\sqrt{n}}
                   = \frac{503.2 - 500}{8/\sqrt{25}}
                   = \frac{3.2}{1.6}
                   = 2.00.
\end{equation}

\textbf{步骤 3：计算 $p$ 值}

因为是双侧检验，$p$ 值为观察到 $|Z| \ge 2.00$ 的概率：
\begin{equation}
    p = 2 \times P(Z \ge 2.00 \mid H_0)
      = 2 \times \bigl(1 - \Phi(2.00)\bigr).
\end{equation}

查标准正态分布表，$\Phi(2.00) = 0.9772$：
\begin{equation}
    p = 2 \times (1 - 0.9772) = 2 \times 0.0228 = 0.0456.
\end{equation}

\medskip
\begin{center}
\begin{tikzpicture}[scale=1.1]
    % Background shade for rejection regions
    \fill[red!15] (-4,0) rectangle (-1.96, 2.2);
    \fill[red!15] (1.96,0) rectangle (4, 2.2);
    % Fill p-value area (observed)
    \fill[blue!25] (2.00,0) rectangle (4, 1.1);
    \fill[blue!25] (-4,0) rectangle (-2.00, 1.1);
    % Axes
    \draw[-stealth] (-4,0) -- (4,0) node[right] {$Z$};
    \draw[-stealth] (0,0) -- (0,2.5) node[above] {密度};
    % Normal curve
    \draw[domain=-3.8:3.8, smooth, thick] plot (\x, {2*exp(-\x*\x/2)/sqrt(2*pi)});
    % Critical values
    \draw[dashed, red, thick] (-1.96,0) -- (-1.96,{2*exp(-1.96*1.96/2)/sqrt(2*pi)})
        node[above, red] {$-1.96$};
    \draw[dashed, red, thick] (1.96,0) -- (1.96,{2*exp(-1.96*1.96/2)/sqrt(2*pi)})
        node[above, red] {$1.96$};
    % Observed Z
    \draw[dashed, blue, thick] (2.00,0) -- (2.00,{2*exp(-2*2/2)/sqrt(2*pi)})
        node[above right, blue] {$Z_{\text{obs}}\!=\!2.00$};
    \draw[dashed, blue, thick] (-2.00,0) -- (-2.00,{2*exp(-2*2/2)/sqrt(2*pi)});
    % Annotations
    \node[blue] at (2.9, 0.25) {$p/2$};
    \node[blue] at (-2.9, 0.25) {$p/2$};
    \node[red] at (-2.6, 1.7) {拒绝域};
    \node[red] at (2.6, 1.7) {拒绝域};
\end{tikzpicture}
\end{center}

\textbf{步骤 4：做出决策}

\begin{itemize}
    \item $p = 0.0456 < \alpha = 0.05$，拒绝 $H_0$。
    \item 按 $p$ 值解释标准，$0.01 < p \le 0.05$，属于\textbf{中等证据}。
\end{itemize}

\textbf{结论：} 在 $\alpha = 0.05$ 水平下，有显著证据表明灌装量偏离了 500 ml。
样本均值 503.2 ml 比标称值高出 3.2 ml，但 $p = 0.0456$ 接近 0.05，
证据不算极强，建议增加样本量再次检验。

\medskip
\textbf{对比 —— 传统拒绝域法：}

若用拒绝域法，$z_{0.025} = 1.96$，$|Z_{\text{obs}}| = 2.00 > 1.96$，
落入拒绝域，同样拒绝 $H_0$。两种方法结论一致。

但 $p$ 值法额外告诉我们：若 $\alpha$ 设为 0.01（$z_{0.005} = 2.576$），
则 $|Z_{\text{obs}}| = 2.00 < 2.576$，$p = 0.0456 > 0.01$，不拒绝 $H_0$。
这说明结论\textbf{依赖于所选的 $\alpha$}。$p$ 值让读者可以自行判断，
而不必受制于作者事先选择的 $\alpha$。

\medskip
\textbf{延伸 —— 若为单侧检验：}

如果质检员只关心灌装量是否\textbf{偏多}（右侧检验 $H_1: \mu > 500$）：
\begin{equation}
    p_{\text{右侧}} = P(Z \ge 2.00) = 1 - \Phi(2.00) = 0.0228.
\end{equation}
$p = 0.0228 < 0.05$，同样拒绝 $H_0$（证据更强，因为 $p$ 比双侧小一半）。

如果只关心是否\textbf{偏少}（左侧检验 $H_1: \mu < 500$）：
\begin{equation}
    p_{\text{左侧}} = P(Z \le 2.00) = \Phi(2.00) = 0.9772.
\end{equation}
$p = 0.9772 \gg 0.05$，毫无证据拒绝 $H_0$。这与样本均值 $503.2 > 500$ 的方向一致。
\end{example}

% ============================================================
%  小结
% ============================================================
\section*{本章小结}

本章系统介绍了统计推断的另一核心内容——假设检验：

\begin{enumerate}
    \item \textbf{基本框架}：原假设/备择假设的设定、两类错误及其概率 $\alpha, \beta$、拒绝域的构造；
    \item \textbf{正态总体均值检验}：$Z$ 检验（方差已知）、$t$ 检验（方差未知）、两样本与配对检验；
    \item \textbf{正态总体方差检验}：$\chi^2$ 检验（单总体）、$F$ 检验（两总体方差比）；
    \item \textbf{置信区间与检验的对偶性}：置信区间可视为所有不被拒绝的参数值集合；
    \item \textbf{样本容量选取}：根据效应量、$\alpha$、功效等在抽样前确定所需 $n$；
    \item \textbf{分布拟合检验}：Pearson $\chi^2$ 检验判断总体是否服从指定分布；
    \item \textbf{秩和检验}：不依赖正态假设的非参数方法（Wilcoxon 秩和检验、符号秩检验）；
    \item \textbf{$p$ 值检验法}：提供证据强度的连续度量，是现代统计实践的标准。
\end{enumerate}

假设检验与参数估计共同构成了\textbf{统计推断}的两大基石。在机器学习中，假设检验广泛应用于 A/B 测试、特征选择、模型比较与分布漂移检测等场景。

% ============================================================
%  第 9 章：方差分析及回归分析
% ============================================================
